\clearpage
\SourceLeaf{357}\pdfbookmark[1]{Nachlass und Veröffentlichung}{parabolisch-veroeffentlichung}
\begin{center}
{\large NACHLASS UND VERÖFFENTLICHUNG.\par}
\bigskip
\pdfbookmark[2]{II. Zur Berechnung der wahren Anomalie in der parabolischen Bewegung}{parabolisch-anomalie}
\textbf{[II. Zur Berechnung der wahren Anomalie in der parabolischen Bewegung.]}\par
\bigskip
\pdfbookmark[3]{1. Tabula nova motus parabolici}{tabula-nova}
\textbf{[1. Tabula nova motus parabolici.]}\par
\medskip
\textbf{Explicatio tabulae novae motus parabolici.}
\end{center}
\medskip
\pdfbookmark[4]{1.}{tabula-nova-001}
\begin{center}1.\end{center}
\noindent\hspace*{\parindent}Denotabimus, perinde ut in Theoria Motus art. 18–20, distantiam a Sole in perihelio per \(q\), anomaliam veram per \(v\), tempus inde a transitu per perihelium per \(t\), denique per [\(k\)] constantem \(3548''{,}188\) (art. [6]). Quum in motu parabolico cometae, cuius massam negligimus, velocitas angularis in secundis expressa fiat generaliter
\[
\frac{\mathrm{d}v}{\mathrm{d}t}=\frac{k\BareRoot{2}}{q^{\frac32}}\cos\tfrac12v^4,
\]
statuemus brevitatis caussa
\[
\frac{k\BareRoot{2}}{q^{\frac32}}=A,
\]
ita ut \(A\) exhibeat velocitatem angularem in perihelio, quam elementis parabolicis cuiusvis cometae adiungere conveniet. Est vero
\[
\log k\BareRoot{2}=3{,}700\ 5216.
\]
\par\medskip
\pdfbookmark[4]{2.}{tabula-nova-002}
\begin{center}2.\end{center}
\noindent\hspace*{\parindent}Relatio inter anomaliam veram atque tempus exprimitur per formulam
\SourceLeaf{358}
\[
\operatorname{tang}\tfrac12v+\tfrac13\operatorname{tang}\tfrac12v^3
=\frac{At}{2\mathpunct{.}206265''}=\frac{\pi}{1296000''}At
\]
denotante \(\pi\) semicircumferentiam circuli cuius radius est unitas. Habetur itaque tempus per anomaliam veram directe; pro calculo inverso, \textsc{Cardan}i regula, introductis duobus angulis auxiliaribus \(P\), \(Q\), formulas sequentes suppeditat:
\[
\operatorname{tang}P=\frac{4\mathpunct{.}206265''}{3At}=\frac{864000''}{\pi At}
\]
\[
\operatorname{tang}Q=\BareRoot[3]{\operatorname{tang}\tfrac12P}
\]
\[
\operatorname{tang}\tfrac12v=2\operatorname{cotang}2Q.
\]
\par\medskip
\pdfbookmark[4]{3.}{tabula-nova-003}
\begin{center}3.\end{center}
\noindent\hspace*{\parindent}At quoniam problema posterius frequentissime occurrit, taediosum foret, semper ad has formulas vel ad methodos indirectas confugere; quamobrem astronomi ad hunc finem tabula auxiliari uti solent. Tabula vulgaris pro argumento
\[
\frac{t}{q^{\frac32}}=\frac{At}{k\BareRoot{2}}
\]
anomaliam veram suppeditat. Cel. \textsc{Burckhardt} illius quantitatis logarithmum adoptavit\footnote{Connaiss. des tems 1816.}. Tabulae \textsc{Barker}ianae argumentum est anomalia vera, cui valor respondens quantitatis
\[
\frac{k\pi\BareRoot{2}\mathpunct{.}t}{17280''\mathpunct{.}q^{\frac32}}=\frac{\pi At}{17280''}
\]
appositus est. In tabula nova huic operi annexa (quae iam constructa erat, quando tabula cel. \textsc{Burckhardt} promulgata est) logarithmum quantitatis \(At\) pro argumento accepimus, per singulas partes millesimas unitatis progredientem. Quum in motu parabolico valor ipsius \(At\) a \(0\) in infinitum progrediatur, hac ratione tabula manifesto neque ab initio neque in fine completa evadere potuisset, quapropter in tabulae initio quantitatem \(At\) ipsam tamquam argumentum accepimus per centenas unitates progredientem usque ad valorem \(40000\) cui respondet anomalia vera \(11^0\ 2'\ 32''\). Dein logarithmi progrediuntur a valore \(4{,}6\) usque ad \(8{,}2\), quibus limitibus respondent anomaliae \(10^0\ 59'\ 26''\) atque \(169^0\ 0'\ 15''\).
\par
\SourceLeaf{359}
Rarissime anomaliis maioribus opus erit. Attamen ut tabula ab utraque parte completa evadat, adiecimus partem tertiam, in qua pro argumento accepimus quantitatem
\[
\BareRoot[3]{\frac{B}{t}}\mathpunct{\raisebox{.45ex}{,}}
\]
ubi
\[
B=\frac{1296000^4}{3\pi^4A}=\frac{1296000^4}{3\pi^4k\BareRoot{2}}q^{\frac32}.
\]
\par
Hoc argumentum decrescendo a valore \(40000\) usque ad \(0\) extenditur, quibus respondent anomaliae \(168^0\ 49'\ 9''\) atque \(180^0\). Ad inveniendum logarithmum constantis \(B\), adscribimus logarithmos sequentes
\[
\log\frac{1296000^4}{3\pi^4}=[22{,}461\ 8207]
\]
\[
\log\frac{1296000^4}{3\pi^4k\BareRoot{2}}=[18{,}761\ 2991.]
\]
\par
Usus partis primae tabulae itaque extenditur
\[
\begin{array}{rll}
\text{a}&t=0&\text{usque ad }[t=]\ 7{,}971\,q^{\frac32}\\
\text{partis secundae a}&t=7{,}934\,q^{\frac32}&\text{usque ad }t=31585\,q^{\frac32}\\
\text{partis tertiae a}&t=30061\,q^{\frac32}&\text{in infinitum.}
\end{array}
\]
\par
Ceterum observamus tabulam integram ad partes centies millesimas minutorum secundorum anomaliae ab initio constructam esse, ut figurae secundae semper tuto fidere liceret.
\par\medskip
\pdfbookmark[4]{4.}{tabula-nova-004}
\begin{center}4.\end{center}
\noindent\hspace*{\parindent}Operae pretium erit, tabulae nostrae initium et finem seorsim considerare, atque anomaliam in seriem explicare secundum potestates argumenti progredientem.
\par
Statuamus
\[
\operatorname{tang}\tfrac12v=\theta
\]
\[
\frac{At}{206265''}=\alpha
\]
eritque
\[
\theta=\tfrac12\alpha-\tfrac13\theta^3.
\]
Hinc eruitur adiumento theorematis \textsc{Lagrang}iani generaliter
\SourceLeaf{360}
\[
\begin{aligned}
\theta^n={}&\left(\frac{\alpha}{2}\right)^n-\frac{n}{3}\left(\frac{\alpha}{2}\right)^{n+2}
+\frac{n\mathpunct{.}n+5}{3\mathpunct{.}6}\left(\frac{\alpha}{2}\right)^{n+4}\\
&-\frac{n\mathpunct{.}+n+7\mathpunct{.}n+8}{3\mathpunct{.}6\mathpunct{.}9}\left(\frac{\alpha}{2}\right)^{n+6}\\
&+\frac{n\mathpunct{.}n+9\mathpunct{.}n+10\mathpunct{.}n+11}{3\mathpunct{.}6\mathpunct{.}9\mathpunct{.}12}\left(\frac{\alpha}{2}\right)^{n+8}-\text{etc.}
\end{aligned}
\]
Substituendo hanc expressionem singularum potestatum \(\theta^n\) in serie
\[
v=2\theta-\tfrac23\theta^3+\tfrac25\theta^5-\tfrac27\theta^7+\text{etc.}
\]
habemus
\[
\begin{aligned}
v={}&\alpha-\tfrac13\cdot\tfrac14\left(1+\tfrac33\right)\alpha^3\\
&+\tfrac15\cdot\tfrac1{16}\left(1+\tfrac53+\tfrac{5\mathpunct{.}6}{3\mathpunct{.}6}\right)\alpha^5\\
&-\tfrac17\cdot\tfrac1{64}\left(1+\tfrac73+\tfrac{7\mathpunct{.}8}{3\mathpunct{.}6}+\tfrac{7\mathpunct{.}8\mathpunct{.}9}{3\mathpunct{.}6\mathpunct{.}9}\right)\alpha^7\\
&+\tfrac19\cdot\tfrac1{256}\Bigl(1+\tfrac93+\tfrac{9\mathpunct{.}10}{3\mathpunct{.}6}+\tfrac{9\mathpunct{.}10\mathpunct{.}11}{3\mathpunct{.}6\mathpunct{.}9}\\
&\hspace{5em}+\tfrac{9\mathpunct{.}10\mathpunct{.}11\mathpunct{.}12}{3\mathpunct{.}6\mathpunct{.}9\mathpunct{.}12}\Bigr)\alpha^9-\text{etc.,}
\end{aligned}
\]
cuius seriei lex est obvia. Eruendo coëfficientium valores numericos fit
\[
\begin{aligned}
v={}&\alpha-\tfrac16\alpha^3+\tfrac{13}{240}\alpha^5-\tfrac{43}{2016}\alpha^7\\
&+\tfrac{191}{20736}\alpha^9-\tfrac{479}{114048}\alpha^{11}+\text{etc.}
\end{aligned}
\]
\par
Manifesto hic \(v\) tamquam arcus in partibus radii expressus considerandus est. Qui si in minutis secundis desideratur, series per \(206265''\) multiplicanda est, ita ut fiat
\[
v=At-\tfrac16\frac{(At)^3}{206265^2}+\tfrac{13}{240}\frac{(At)^5}{206265^4}-\text{etc.}
\]
Anomaliae parvae itaque ab argumentis respondentibus parum differunt.
\par
Ceterum expressio pro \(\theta^n\) supra data etiam seriem elegantem pro radio vectore \(r\) suppeditat; fit enim
\[
\begin{aligned}
r={}&q(1+\theta\theta)\\
={}&q\Bigl(1+\tfrac14\alpha\alpha-\tfrac13\cdot\frac{\alpha^4}{8}+\frac{7}{3\mathpunct{.}6}\frac{\alpha^6}{32}\\
&-\frac{9\mathpunct{.}10}{3\mathpunct{.}6\mathpunct{.}9}\frac{\alpha^8}{128}
+\frac{11\mathpunct{.}12\mathpunct{.}13}{3\mathpunct{.}6\mathpunct{.}9\mathpunct{.}12}\frac{\alpha^{10}}{512}\\
&-\frac{13\mathpunct{.}14\mathpunct{.}15\mathpunct{.}16}{3\mathpunct{.}6\mathpunct{.}9\mathpunct{.}12\mathpunct{.}15}\frac{\alpha^{12}}{2048}+\text{etc.}\Bigr).
\end{aligned}
\]
Ostendi potest utramque seriem convergere quamdiu \(\alpha<\tfrac43\) sive \(v<61^0\ 35'\ 45''{,}4\).
\par
